\chapter{Evaluation}
\label{ch:evaluation}

Dieses Kapitel wertet die Roundtrip-Latenz des Tunnels im Docker-Testbett unter
emulierten Netzbedingungen aus und beantwortet damit die Forschungsfragen FF2
(Latenz-Overhead) und FF3 (Verhalten unter Netzbedingungen).

\section{Testbett und Methodik}
\label{sec:testbett}

Das Testbett ist vollständig containerisiert (Docker Compose): Edge, Agent,
Origin (ein TCP-Echo-Dienst) und Client laufen auf einem statischen Netz. Die
Link-Impairment wird mit \texttt{tc netem} auf der Egress-Schnittstelle des Edge
gesetzt und ist über Umgebungsvariablen (Verzögerung, Verlust, Bandbreite)
parametrisiert. Der Client misst im Bench-Modus je Bedingung $n=5$ vollständige
Roundtrips -- Verbindungsaufbau, PoW-Rendezvous, Tunnel und Echo-Verifikation --
jeweils mit \emph{frisch aufgebauter} QUIC-Verbindung.

Die Auswertung ist vollständig reproduzierbar und ihrerseits containerisiert:
\texttt{scripts/sweep.sh} fährt die Bedingungsmatrix ab und sammelt die Ergebnisse
als CSV, \texttt{scripts/plot.sh} erzeugt die Abbildungen (Matplotlib in einem
Python-Container) und \texttt{scripts/tabulate.py} die Tabelle. Vermessen wurde
die Matrix Verzögerung $\in \{0, 20, 50\}$\,ms $\times$ Verlust $\in \{0, 2\}$\,\%.

\section{Ergebnisse}
\label{sec:ergebnisse}

Tabelle~\ref{tab:latency} fasst alle Messpunkte zusammen;
Abbildung~\ref{fig:latency-delay} zeigt die Latenz über der Verzögerung,
Abbildung~\ref{fig:latency-loss} über dem Paketverlust.

% Ergebnistabelle (generiert, scripts/tabulate.py).
\begin{table}[t]
  \centering
  \caption{Tunnel-Roundtrip-Latenz unter emulierten Netzbedingungen (\texttt{tc netem}).}
  \label{tab:latency}
  \begin{tabular}{llrrrrrr}
    \toprule
    Verzögerung & Verlust & $n$ & Mittel & p50 & p95 & Min & Max \\
     &  &  & (ms) & (ms) & (ms) & (ms) & (ms) \\
    \midrule
    0ms & 0\% & 5 & 6.8 & 6.6 & 7.9 & 5.8 & 7.9 \\
    0ms & 2\% & 5 & 6.9 & 6.9 & 8.2 & 5.5 & 8.2 \\
    20ms & 0\% & 5 & 92.7 & 88.7 & 109.2 & 88.1 & 109.2 \\
    20ms & 2\% & 5 & 92.2 & 88.2 & 108.5 & 88.0 & 108.5 \\
    50ms & 0\% & 5 & 217.9 & 207.5 & 260.4 & 205.5 & 260.4 \\
    50ms & 2\% & 5 & 237.5 & 207.7 & 306.5 & 207.5 & 306.5 \\
    \bottomrule
  \end{tabular}
\end{table}


\begin{figure}[t]
  \centering
  \includegraphics[width=0.8\linewidth]{latency-vs-delay.png}
  \caption{Tunnel-Roundtrip über der Edge-Verzögerung, je Verlustwert eine
    Serie; der obere Fehlerbalken markiert das p95.}
  \label{fig:latency-delay}
\end{figure}

\begin{figure}[t]
  \centering
  \includegraphics[width=0.8\linewidth]{latency-vs-loss.png}
  \caption{Tunnel-Roundtrip über dem Paketverlust, je Verzögerungswert eine
    Serie.}
  \label{fig:latency-loss}
\end{figure}

\section{Diskussion}
\label{sec:diskussion}

\paragraph{FF2 -- Latenz-Overhead.} Ohne Netzverzögerung beträgt der
Roundtrip im Mittel 6{,}8\,ms; dies ist der Eigenanteil des Tunnels aus
QUIC-Handshake, PoW-Rendezvous und Mehrsprung-Relay. Über die Messpunkte
(0\,ms~$\rightarrow$~6{,}8; 20\,ms~$\rightarrow$~92{,}7; 50\,ms~$\rightarrow$~217{,}9)
skaliert die Latenz näherungsweise linear:
$\text{RT} \approx 6{,}8 + 4{,}2 \cdot d$ mit $d$ = einseitige Edge-Verzögerung
in Millisekunden. Der Faktor~$\approx$~4{,}2 spiegelt wider, dass jeder
Anwendungs-Roundtrip die verzögerte Egress mehrfach quert (Verbindungsaufbau,
Rendezvous und der Relay-Pfad Client$\rightarrow$Edge$\rightarrow$Agent$\rightarrow$Origin
und zurück). Da pro Iteration frisch verbunden wird, trägt der Handshake einen
festen Vielfachen der Umlaufzeit bei.

\paragraph{FF3 -- Verhalten unter Paketverlust.} Bei niedriger Umlaufzeit
(0\,ms, 20\,ms) ist ein Verlust von 2\,\% im Mittel nicht messbar
(6{,}8$\rightarrow$6{,}9\,ms bzw.\ 92{,}7$\rightarrow$92{,}2\,ms, innerhalb der
Streuung). Bei 50\,ms hingegen steigt der Mittelwert um rund 9\,\%
(217{,}9$\rightarrow$237{,}5\,ms) und das p95 um rund 18\,\%
(260{,}4$\rightarrow$306{,}5\,ms), während der Median nahezu unverändert bleibt
(p50~$\approx$~207\,ms). Verlustbedingte Neuübertragungen wirken sich also
proportional zur Umlaufzeit aus und zeigen sich primär im Tail (p95), nicht im
Median.

\section{Limitierungen}
\label{sec:limitierungen}

Die Ergebnisse sind unter drei Einschränkungen zu lesen. Erstens wirkt die
Impairment nur auf der Edge-Egress und damit nicht symmetrisch; die absolute
Latenz ist eine untere Schranke für eine voll bidirektional verzögerte Strecke.
Zweitens dient $n=5$ hier der Demonstration der reproduzierbaren Mess-Pipeline,
nicht einer statistisch belastbaren Aussage -- größere Stichproben und eine
breitere Matrix (inklusive Bandbreitenbegrenzung) sind über dieselben Skripte
unmittelbar zugänglich. Drittens überzeichnet der Frisch-Verbindungsaufbau je
Iteration den Handshake-Anteil; ein Keep-Alive-Datenpfad läge niedriger. Diese
Punkte greift der Ausblick (Kapitel~\ref{ch:fazit}) auf.
